\section{Introduction}
\label{sec:intro}

Reinforcement learning with verifiable rewards (RLVR) makes a language model
solve more problems on the first try and fewer problems when given many tries:
it raises pass@1 while the base model catches up and overtakes it at large
sampling budget $k$ \citep{yue2025rlincentivize}. Whether this means RLVR fails to
add reasoning, or only needs more training to pull ahead \citep{liu2025prorl},
has been argued from curves. We settle it with a theorem. The reversal is
neither a failure nor a transient: it is a \emph{price}, and we compute the
exchange rate at which it is paid.

We model RLVR as what it is in the population limit, an exponential tilt of a
base distribution over solution strategies that are shared across prompts. One
identity then drives the two phenomena that have been studied apart. On the
strategy simplex, probability is conserved, so reinforcement spends it. When
strategies specialize, each is the only route to some prompts, and the spending
is visible at the level of prompts. This yields a sharp dichotomy. \emph{Every}
update that raises pass@1 lowers the success of some prompt (Theorem
\ref{thm:nfl}); there is no free lunch, and a theorem of the alternative gives
the exact condition, a strictly positive price on prompts that equalizes strategy
fitness. \emph{When} that loss reaches the hardest,
finely covered prompts, it forces the pass@$k$ crossover (Theorem
\ref{thm:price}), and the number of times the curves can cross is bounded by the
sign changes of the shift in per-prompt success, a totally positive argument
that is tight. The same tilt yields entropy collapse and the
entropy-performance law with an identifiable exponent (Theorem
\ref{thm:tiltlaw}). In particular, no update both raises pass@1 and dominates the
base pass@$k$ curve at every $k$: reinforcement cannot widen the reasoning
frontier, it can only move along it. That is the exact sense in which RLVR does
not incentivize reasoning beyond the base model, and the sense in which, under a
single named exception, it can.

The price is not a bug to be patched. A constrained update that forbids any
prompt from regressing removes the crossover exactly, and under specialization
it also removes the pass@1 gain: the only way to pay nothing is to gain nothing.
Between the extremes lies a frontier (Theorem \ref{thm:frontier}), the exact
exchange rate between pass@1 and pass@$k$, which an optimal per-prompt floor
traces out. A diversity bonus helps both metrics only when the unmodified update
was off this frontier, which is what the frontier predicts.

\paragraph{Contributions.}
\begin{itemize}
\item \textbf{No free lunch (Theorem \ref{thm:nfl}).} Under exclusive (or dominant) specialization,
every pass@1-increasing update lowers some prompt's success. We give the exact
Stiemke condition and show the hypothesis is sharp: unique-best is insufficient
(a worked counterexample) and a dominating strategy restores a
free lunch.
\item \textbf{The crossing bound (Theorem \ref{thm:price}a,b).} For any
base-to-RL shift, the number of pass@$k$ crossings is at most the sign variation
of the success-CDF difference; the bound is tight. This is the general statement
that concurrent work leaves open for the kernel case.
\item \textbf{Forced crossover and its scope (Theorem \ref{thm:price}c--e).}
Under a named fine-coverage condition the crossover is forced; without it there
is a per-prompt loss but no crossover. The order-preserving case recovers the
known single-crossover theorem as a corollary.
\item \textbf{The price frontier (Theorem \ref{thm:frontier}).} The optimal
pass@1 versus pass@$k$ exchange rate, with a closed form for two strategies.
\item \textbf{A diagnostic.} A tool that takes only per-prompt base and RL success
rates, available in any evaluation log, and returns the bound, the
order-preserving test, and the price. We read published curves through it and
find the predicted structure.
\end{itemize}

The claims are stated at the level of the model and proved in the appendix; every
number is produced and checked by an accompanying package. We do not claim
anything about a specific language model beyond what published curves show. The
message is one line: in RLVR there is no free lunch, only a price, and we give
its rate.